\chapter{Nasal Obstruction}

\section*{Definition}
Nasal obstruction is the subjective sensation of reduced airflow through one or both nasal passages, which may be constant or alternating (nasal cycle), and can significantly impair quality of life and sleep.

\section*{What Happens in Your Body}
Obstruction arises from structural narrowing of the nasal airway (septal deviation, turbinate hypertrophy, nasal valve collapse) or from mucosal congestion secondary to inflammation (allergic, infectious, or vasomotor). The nasal valve (internal nasal valve angle normally 10--15\textdegree) is the narrowest segment and the most common site of obstruction.

\section*{Causes}
\begin{itemize}
  \item \textbf{Structural:} Nasal septal deviation, septal perforation, nasal valve collapse, turbinate hypertrophy
  \item \textbf{Inflammatory:} Allergic rhinitis, non-allergic (vasomotor) rhinitis, rhinitis medicamentosa (decongestant overuse)
  \item \textbf{Infectious:} Acute or chronic rhinosinusitis, nasal vestibulitis
  \item \textbf{Neoplastic:} Inverted papilloma, juvenile nasopharyngeal angiofibroma, squamous cell carcinoma (rare)
  \item \textbf{Iatrogenic:} Post-surgical synechiae, nasal packing, or foreign body
\end{itemize}

\section*{When to See a Doctor}
See your doctor if:
\begin{itemize}
  \item Unilateral, progressive nasal obstruction (structural or neoplastic cause)
  \item Obstruction associated with epistaxis, facial pain, or anosmia
  \item Snoring, gasping, or witnessed apneas suggesting OSA
  \item Failure of medical therapy (topical steroids, antihistamines) after 4--6 weeks
  \item Congenital nasal obstruction in an infant (suspect choanal atresia)
\end{itemize}

\section*{Related Symptoms}
\begin{itemize}
  \item Mouth breathing and dry mouth upon waking
  \item Snoring and sleep disturbance
  \item Hyponasal speech
  \item Anosmia or hyposmia (reduced sense of smell)
  \item Postnasal drip, rhinorrhea, and sneezing (if inflammatory)
\end{itemize}
\textbf{Important:} Unilateral nasal obstruction with purulent discharge in a child suggests a nasal foreign body until proven otherwise.

\section*{Self-Care / Home Management}
\begin{itemize}
  \item Saline nasal irrigation (Neti pot or squeeze bottle) once or twice daily
  \item Topical decongestant sprays (oxymetazoline) for no more than 3 consecutive days to avoid rhinitis medicamentosa
  \item Intranasal corticosteroid sprays (fluticasone, mometasone) for persistent allergic\/inflammatory obstruction
  \item Nasal dilator strips (external Breathe Right strips) for nocturnal nasal valve collapse
  \item Humidifier in the bedroom to prevent nasal drying
  \item Elevate head of bed to reduce nasal congestion
\end{itemize}

\section*{How Your Doctor Will Evaluate You}
\begin{itemize}
  \item Anterior rhinoscopy and nasal endoscopy (rigid or flexible) to assess the nasal cavity, septum, turbinates, and posterior choanae
  \item Cottle maneuver (lateral traction on the cheek) to assess internal nasal valve collapse
  \item Allergy testing (skin prick or serum specific IgE) if allergic rhinitis suspected
  \item CT sinuses without contrast for chronic obstruction with sinusitis or suspected neoplasm
  \item Nasal peak inspiratory flow or rhinomanometry for objective airflow quantification
  \item Sleep study (polysomnography) if OSA suspected
\end{itemize}

\section*{Treatment Options}
\begin{itemize}
  \item \textbf{Medical:} Intranasal corticosteroids (first-line), oral antihistamines, leukotriene receptor antagonists, allergen immunotherapy (for allergic rhinitis)
  \item \textbf{Turbinate reduction:} Radiofrequency ablation, submucous resection, or coblation for hypertrophic inferior turbinates
  \item \textbf{Septoplasty:} Surgical correction of significant septal deviation
  \item \textbf{Nasal valve repair:} Spreader grafts, flaring sutures, or lateral crus repositioning for nasal valve collapse
  \item \textbf{Functional rhinoplasty:} Combined septal, turbinate, and valve procedures
  \item \textbf{Endoscopic sinus surgery:} For chronic rhinosinusitis-related obstruction
  \item \textbf{Management of rhinitis medicamentosa:} Discontinue decongestant spray, short course of oral steroids, and intranasal steroid spray
\end{itemize}

\section*{References}
\begin{thebibliography}{99}
\bibitem{nasal_obstruction:ref1} Lal D, Corey JP. ``Nasal obstruction: a comprehensive review.'' \textit{Otolaryngologic Clinics of North America}, 2020;53(5):769--780.
\bibitem{nasal_obstruction:ref2} Chandra RK, Patadia MO, Raviv J. ``Diagnosis and management of nasal obstruction.'' \textit{Journal of the American Medical Association}, 2018;320(17):1783--1794.
\bibitem{nasal_obstruction:ref3} Rhee JS, Book DT, Burzynski M, et al. ``Quality of life assessment in nasal airway obstruction.'' \textit{Journal of the American Medical Association Facial Plastic Surgery}, 2003;5(4):327--331.
\bibitem{nasal_obstruction:ref4} Kimmelman CP. ``The problem of nasal obstruction.'' \textit{Otolaryngologic Clinics of North America}, 1989;22(2):253--264.
\bibitem{nasal_obstruction:ref5} Seiden AM. ``Nasal obstruction and the nasal cycle.'' \textit{Otolaryngologic Clinics of North America}, 2005;38(6):1167--1178.
\bibitem{nasal_obstruction:ref6} D'Souza N, Abrol R, Golhar K. ``Management of nasal obstruction.'' \textit{British Medical Journal}, 2018;363:k4239.

\end{thebibliography}
